\chapter{Literature Review}
\label{ch:litreview}

\section{Large Language Models as Recommender Systems}

Recommender systems have traditionally relied on collaborative filtering, content-based filtering, or hybrid approaches, each requiring a structured representation of users and items (Pazzani and Billsus, 2007). The arrival of general-purpose LLMs has prompted a substantial and rapidly growing body of work exploring whether these models can perform recommendation directly, without task-specific training. Liu et al. (2024) and Wu et al. (2024) both survey this emerging area, characterising LLM-based recommendation approaches along a spectrum from using an LLM purely as a feature encoder within a conventional pipeline, through to using the LLM as the recommender itself via prompting --- the approach adopted in this study.

A recurring theme in this literature is that LLMs bring capabilities a conventional recommender lacks: they can incorporate free-text context, generalise to items or users outside the training distribution (zero-shot generalisation), and produce natural-language justifications for their choices (Liu et al., 2024). These properties are directly relevant to a domain like credit-card recommendation, where a user's stated goals (``I want to build credit,'' ``I'm planning a large purchase'') are naturally expressed in language rather than as a fixed numeric feature. At the same time, the same surveys consistently flag three risks that motivate the evaluation framework adopted in this dissertation: (i) sensitivity to how the task is phrased in the prompt, (ii) hallucination of facts not present in the supplied context, and (iii) a lack of the determinism and auditability that a conventional, formula-based recommender guarantees by construction.

\section{Prompt Engineering and Prompting Strategies}

Since Brown et al. (2020) demonstrated that GPT-3 could perform new tasks from a handful of in-context examples without any parameter updates --- ``few-shot learning'' --- prompt design has become a first-class variable in how LLMs are applied to downstream tasks, rather than a fixed implementation detail. Schulhoff et al. (2024), in one of the most comprehensive taxonomies of prompting techniques to date, catalogue dozens of distinct prompting techniques for text-based tasks alone, underlining that ``zero-shot'' and ``few-shot'' are two points on a much larger design space that also includes explicit output-format constraints (referred to here as ``structured'' prompting).

Empirical comparisons of zero-shot against few-shot prompting report task-dependent results. Few-shot prompting has been shown to yield meaningful accuracy gains over zero-shot in several NLP classification tasks, with in-context examples providing the model with an implicit template for both reasoning and expected output format, and gains are typically achieved with as few as one to five examples. However, this advantage is not universal: for some domains, particularly complex reasoning and specialised tasks, well-specified zero-shot instructions have been found to be competitive with, or in some cases superior to, few-shot prompting, suggesting that the relative benefit of providing examples depends heavily on how much the task benefits from an explicit template versus additional worked examples. This inconsistency in the literature is itself a motivation for this dissertation's comparative design: rather than assuming \textit{a priori} that one strategy dominates, the three strategies are evaluated empirically against a common baseline and a common evaluation framework.

The ``structured'' prompting condition used in this dissertation --- an explicit output schema (JSON, with fields for card ID, name, and justification) specified in the prompt --- reflects a widely used mitigation for a separate, practical problem: unconstrained natural-language output is difficult to parse reliably. Requiring the model to populate a fixed schema is one of the more common patterns discussed in the prompting-technique literature (Schulhoff et al., 2024) precisely because it improves downstream parseability, which matters for an evaluation pipeline like the one used here that must automatically extract and score card recommendations from raw model output at scale (900 responses).

\section{Content-Based Filtering as an Evaluation Baseline}

Content-based recommendation constructs a profile of item attributes and a corresponding profile of user preferences over those same attributes, then recommends items whose profile most closely matches the user's, typically via cosine similarity or a related distance metric (Pazzani and Billsus, 2007; Lops, de Gemmis and Semeraro, 2011). This family of methods is well established, fully deterministic, and --- critically for this dissertation's purposes --- entirely independent of any language model, making it an appropriate conventional comparator for LLM-generated recommendations. This dissertation's baseline (Section~\ref{sec:baseline}) follows this established approach directly: each of the 20 catalogue cards and each of the 300 synthetic profiles are encoded as numeric feature vectors over the same set of attributes (annual fee, category-specific reward rates, foreign-transaction fee), and cosine similarity between a profile vector and every card vector determines the top-3 reference recommendation for that profile.

\section{Evaluating LLM Output: Quality, Consistency, and Factuality}

Evaluating LLM-generated text is itself an active research area, distinct from evaluating a classification or ranking model's accuracy, because LLM output is unstructured, open-ended, and often has no single ``correct'' answer to compare against. Zheng et al. (2023) introduced the ``LLM-as-a-judge'' paradigm --- using a strong LLM to score the outputs of another model against a rubric --- after showing that GPT-4 judgments agreed with crowdsourced human preferences over 80\% of the time, comparable to human--human agreement, on their MT-Bench and Chatbot Arena benchmarks. They also documented specific failure modes of this approach relevant to any study using it, including position bias, verbosity bias, and self-enhancement bias. This dissertation's evaluation design (Section~\ref{sec:evalframework}) mitigates self-enhancement bias specifically by using a larger, different model (Llama 3.3 70B) to judge the outputs of the smaller model under evaluation (Llama 3.1 8B), rather than having a model judge its own output.

Two further evaluation properties matter specifically for a recommendation task rather than a general chat task. First, factual correctness --- whether specific facts an LLM states about a recommended product are actually true --- can, unlike open-ended explanation quality, be checked deterministically when a fixed, known catalogue of ground-truth facts exists, as is the case here (Section~\ref{sec:evalframework}). Second, consistency --- whether a model gives materially the same answer when queried repeatedly with the same input --- is a property with no equivalent in a deterministic baseline, but is essential to any claim that an LLM-based recommender could be relied upon in practice; a system whose recommendation changes at random between identical requests is difficult to trust regardless of its average accuracy.

\section{Financial Product Recommendation}

Recommendation in financial services has typically relied on the same collaborative- and content-based methods used elsewhere, applied to transaction histories, account data, or stated preferences (Met, 2024). More recent work has explored deep learning and transformer-based approaches for financial product recommendation specifically, citing the need to capture more complex relationships between customer behaviour sequences and product attributes than simple similarity metrics allow. This dissertation sits at the intersection of this financial-recommendation literature and the LLM-recommendation literature reviewed in Section~2.1: it applies an LLM, rather than a purpose-trained deep learning model, to a financial product recommendation task, and evaluates it against the same kind of content-based baseline that has historically been the standard approach in this domain.

\section{Critical Synthesis and Positioning of This Study}

Read together, the four strands of literature reviewed above expose a structural gap rather than a purely empirical one. The LLM-recommendation surveys (Section 2.1) establish \textit{that} LLMs can be used as recommenders and catalogue \textit{how}, but are largely architecture- and technique-focused; they do not, in general, interrogate whether a given prompting strategy is actually reliable enough to trust in a specific, high-stakes consumer domain --- the surveys treat prompting as one design choice among many rather than as a variable worth isolating and testing empirically against a domain-appropriate baseline. Conversely, the prompt-engineering literature (Section 2.2) isolates prompting strategy as its central variable, but the reported zero-shot/few-shot comparisons are drawn almost entirely from general NLP benchmarks --- classification, sentiment, biomedical QA --- not from recommendation tasks with a real product catalogue and a domain-specific correctness criterion. Neither strand, taken alone, answers the question this dissertation asks: for a bounded, factually-checkable recommendation task with real product data, does prompting strategy matter, and does it matter for the same reasons (accuracy) or for different ones (consistency, factual reliability) than in general NLP tasks?

This distinction matters methodologically. Most of the zero-shot/few-shot comparisons surveyed in Section 2.2 use accuracy or F1 against a single correct label as their only outcome measure. A recommendation task has no single correct answer in the same sense --- the baseline in Section 2.3 provides a principled reference point, not a ground truth in the classification sense --- which is precisely why this dissertation's evaluation framework (Section~\ref{sec:evalframework}) deliberately combines four distinct criteria (overlap with baseline, factual correctness, explanation quality, consistency) rather than a single accuracy figure. This is a direct methodological response to a limitation visible across the reviewed prompt-engineering literature: a model can be prompted into a \textit{more accurate-looking} answer while becoming less consistent or less factually grounded, and a single-metric evaluation would not detect that trade-off. The results in Chapter~\ref{ch:analysis} bear this out directly.

The financial-recommendation literature (Section 2.5) and the LLM-as-judge literature (Section 2.4) each fill a different part of the remaining gap: the former supplies the domain and the conventional baseline this dissertation compares against, while the latter supplies a defensible, literature-grounded method for scoring the one evaluation criterion (explanation quality) that cannot be checked deterministically. No source identified in this review combines all three elements --- a real, factually-checkable product catalogue; a controlled, within-subjects comparison of prompting strategies; and a mixed quantitative/automated-qualitative evaluation framework covering accuracy, factuality, and stability together --- which is the specific gap this dissertation addresses.
